\scanpage{106}
% A vertical pen stroke groups the statement at the left.
\uline{Corollaire} Avec les notations précédentes, si
$\underline v\in G^+$, $s\in S=\mathrm{Im}(\underline S\to E)$, alors
\[
\underline v.s=s\Longleftrightarrow v=\mathrm{id}\qquad (v=\varphi_G(\underline v)\in\mathrm{Aff}(E))
\]

\deleted{[\unsure{Mais} b) signifie que $\underline v s_{\unsure{0}}=s_{\unsure{0}}$ implique $\underline v\underline s_0=\underline s_0$]}
Ainsi,
\diagram{\begin{tikzcd}[ampersand replacement=\&,column sep=small,row sep=tiny]
\text{a)}\arrow[dr,Rightarrow]\&{}\&{}\&{}\&{}\\
{}\&\text{b$'$)}\arrow[dr,Rightarrow]\&{}\&{}\&{}\\
{}\&{}\&\text{b)}\arrow[r,Leftrightarrow]\&\text{c)}\arrow[r,Leftrightarrow]\&\text{d)}
\end{tikzcd}}
\deleted{d'ailleurs}

\deleted{et b) $\unsure{-|\Rightarrow}$ b$'$), d'ailleurs \unsure{$\neg$ b)} signifie \deleted{\unsure{implique}}.
$\exists d<n$ tel que \deleted{$u^ds_0=s_0$} \inserted{$u^d=\mathrm{id}$} implique que
$u^da_0=a_0$ i.e. $a_d=a_0$ donc non b$'$), donc}

Il reste à prouver b), c), d) impliquent a).
Soit donc $\underline s\in\underline S$, $\underline a\in\underline A$ supposons $s\in a$
prouvons $\underline s$ incident à $\underline a$. Soient $\underline s_0,\underline s_1$ les
deux sommets \deleted{\unsure{de}} incidents à $\underline a$ \inserted{(distincts)}, donc $s_0,s_1\in a$,
\inserted{et $s_0\ne s_1$,} il faut prouver que \deleted{\unsure{si $s$}} $s=s_0$ ou $s=s_1$,
(ce qui implique $\underline s=\underline s_0$ ou $\underline s=\underline s_1$ donc
$\underline s$ incident à $\underline a$). \deleted{\unsure{Fais|Faux}} Mais si on avait
$s\ne s_0,s_1$, alors $\mathrm{Card}\,a\cap S>2$ (où $S=$
$\mathrm{Im}(\underline S\unsure{-|\to}E)$) a fortiori $\mathrm{Card}\,a\cap C>2$,
où $C$ est la conique circonscrite canonique.
donc $a\subset C$, donc $C$ est \deleted{dégénérée} \inserted{singulière}, donc
$\alpha=-2$. Mais le cas $\alpha=-2$ a été examiné,
et on trouve $C=d\cup d'$, $d$ et $d'$ deux
droites parallèles \uline{distinctes}, avec $s_i\in d$
pour $i$ pair, $s_i\in d'$ pour $i$ impair, donc
$\mathrm{dr}(s_0,s_{\unsure{1|j}})\cap C=\{s_0,s_{\unsure{1|j}}\}$, absurde -- \unsure{gagné}.
% archive sheet number: 105
